% =============================================================================
% CONTINUATION-IN-PART PATENT APPLICATION
% TEMPORAL DIMENSION AUGMENTATION AND MULTI-MODAL MANIFOLD NAVIGATION
% FOR TIMESTAMPED EMBEDDING SPACES
%
% Inventor:  Eric G. Suchanek, PhD
% Assignee:  Flux-Frontiers
% Filing Date: 2026-03-27
%
% This application is a continuation-in-part of co-pending U.S. Patent
% Application [PARENT APPLICATION NUMBER], filed 2026-03-26, entitled
% "System and Method for Offline Semantic Pre-Computation and Knowledge
% Graph Construction from Unstructured Text Corpora."
% =============================================================================

\documentclass[12pt,letterpaper]{article}

\usepackage[
  top=1.0in,
  bottom=1.0in,
  left=1.25in,
  right=1.25in
]{geometry}

\usepackage{setspace}
\usepackage{parskip}
\usepackage{titlesec}
\usepackage{enumitem}
\usepackage{booktabs}
\usepackage{array}
\usepackage{longtable}
\usepackage{amsmath}
\usepackage{amssymb}
\usepackage{listings}
\usepackage{xcolor}
\usepackage{fancyhdr}
\usepackage{lastpage}
\usepackage{hyperref}
\usepackage{microtype}
\usepackage{caption}
\usepackage{float}

% --------------------------------------------------------------------------
% Hyperref configuration (no colored boxes in printed patent)
% --------------------------------------------------------------------------
\hypersetup{
  colorlinks=false,
  pdfborder={0 0 0},
  pdftitle={CIP: Temporal Dimension Augmentation and Multi-Modal Manifold Navigation},
  pdfauthor={Eric G. Suchanek, PhD},
}

% --------------------------------------------------------------------------
% Header / footer
% --------------------------------------------------------------------------
\pagestyle{fancy}
\fancyhf{}
\renewcommand{\headrulewidth}{0.4pt}
\lhead{\small Continuation-in-Part Application}
\chead{\small TEMPORAL DIMENSION AUGMENTATION AND MANIFOLD NAVIGATION}
\rhead{\small Flux-Frontiers}
\cfoot{\small Page \thepage\ of \pageref{LastPage}}

% --------------------------------------------------------------------------
% Section formatting (patent style: bold, no indentation)
% --------------------------------------------------------------------------
\titleformat{\section}{\bfseries\large}{\thesection.}{0.5em}{}
\titleformat{\subsection}{\bfseries\normalsize}{\thesubsection.}{0.5em}{}
\titleformat{\subsubsection}{\bfseries\normalsize\itshape}{\thesubsubsection.}{0.5em}{}

% --------------------------------------------------------------------------
% Listings (Python source code figures)
% --------------------------------------------------------------------------
\definecolor{codebg}{rgb}{0.97,0.97,0.97}
\definecolor{codecomment}{rgb}{0.35,0.55,0.35}
\definecolor{codestring}{rgb}{0.63,0.13,0.13}
\definecolor{codekeyword}{rgb}{0.00,0.00,0.75}
\definecolor{codenumber}{rgb}{0.50,0.50,0.50}

\lstdefinestyle{pythonstyle}{
  language=Python,
  backgroundcolor=\color{codebg},
  basicstyle=\ttfamily\small,
  keywordstyle=\color{codekeyword}\bfseries,
  stringstyle=\color{codestring},
  commentstyle=\color{codecomment}\itshape,
  numberstyle=\tiny\color{codenumber},
  numbers=left,
  numbersep=8pt,
  stepnumber=1,
  frame=single,
  framerule=0.4pt,
  breaklines=true,
  breakatwhitespace=true,
  tabsize=4,
  showspaces=false,
  showstringspaces=false,
  captionpos=b,
  xleftmargin=1.5em,
  xrightmargin=0.5em,
  aboveskip=1em,
  belowskip=0.5em,
}

\lstset{style=pythonstyle}

% --------------------------------------------------------------------------
% Claim numbering environment
% --------------------------------------------------------------------------
\newcounter{claimcounter}
\setcounter{claimcounter}{0}

\newenvironment{claim}{%
  \refstepcounter{claimcounter}%
  \noindent\textbf{Claim \theclaimcounter.}\ %
}{%
  \par\medskip%
}

% --------------------------------------------------------------------------
% Convenience macros
% --------------------------------------------------------------------------
\newcommand{\parentapp}{U.S. Patent Application [PARENT APPLICATION NUMBER]}
\newcommand{\parenttitle}{\emph{System and Method for Offline Semantic
Pre-Computation and Knowledge Graph Construction from Unstructured Text Corpora}}
\newcommand{\turtlend}{\texttt{TurtleND}}
\newcommand{\temporalflyer}{\texttt{TemporalFlyer}}
\newcommand{\augtime}{\texttt{augment\_with\_time}}
\newcommand{\expanddim}{\texttt{expand\_dim()}}
\newcommand{\orientintime}{\texttt{orient\_in\_time()}}
\newcommand{\orienttoward}{\texttt{orient\_toward()}}
\newcommand{\alphabar}{$\alpha$}

% =============================================================================
\begin{document}
% =============================================================================

% --------------------------------------------------------------------------
% Title block
% --------------------------------------------------------------------------
\begin{center}
  {\LARGE\bfseries PATENT APPLICATION}\\[0.5em]
  {\Large\bfseries CONTINUATION-IN-PART}\\[1.5em]
  {\large\bfseries TEMPORAL DIMENSION AUGMENTATION AND MULTI-MODAL MANIFOLD
  NAVIGATION FOR TIMESTAMPED EMBEDDING SPACES}\\[2em]
  \begin{tabular}{ll}
    \textbf{Inventor:}        & Eric G. Suchanek, PhD \\
    \textbf{Assignee:}        & Flux-Frontiers \\
    \textbf{Filing Date:}     & 2026-03-27 \\
    \textbf{Application Type:} & Continuation-in-Part (CIP) \\
    \textbf{Parent Application:} & \parentapp \\
    \textbf{Parent Filed:}    & 2026-03-26 \\
  \end{tabular}
\end{center}

\vspace{1em}
\hrule
\vspace{2em}

\tableofcontents

\newpage

% =============================================================================
\section*{CROSS-REFERENCE TO RELATED APPLICATIONS}
\addcontentsline{toc}{section}{Cross-Reference to Related Applications}
% =============================================================================

This application is a continuation-in-part of \parentapp, filed 2026-03-26,
titled \parenttitle, the entire disclosure of which is incorporated herein by
reference. The parent application discloses and claims a general system and
method for offline semantic pre-computation of unstructured text corpora and
construction of a hybrid knowledge graph therefrom, encompassing corpus
ingestion, multi-phase NLP transformation, metadata extraction and direct write
to structured chunk frontmatter, knowledge graph construction, multi-process
parallel embedding, and JSON embedding cache production. The parent application
applies to any unstructured text corpus and is not limited to timestamped
corpora.

The present continuation-in-part application discloses and claims subject matter
not disclosed in the parent application, specifically: (1) configurable temporal
dimension augmentation of pre-computed embedding vectors; (2) the
\augtime\ algorithm and its parameterization; (3) multi-modal manifold navigation
with time primitives \expanddim, \orientintime, and \orienttoward; (4) three
distinct flight modes (semantic, temporal, mixed) supported by a
\temporalflyer\ navigator; and (5) the generalization of scalar-dimension
augmentation to arbitrary navigable axes beyond time. The present invention
constitutes a specialization of the parent's general pipeline to the class of
corpora possessing temporal metadata: it exploits the timestamp field that the
parent's direct-metadata-write mechanism places in chunk frontmatter and promotes
that field from a post-retrieval filter to a first-class geometric dimension of
the embedding space.

All elements of the parent application's pipeline --- including corpus ingestion,
multi-phase NLP transformation, knowledge graph construction, multi-process
parallel embedding, and JSON embedding cache --- are incorporated herein by
reference and serve as the foundation upon which the present invention operates.

% =============================================================================
\section*{FIELD OF THE INVENTION}
\addcontentsline{toc}{section}{Field of the Invention}
% =============================================================================

The present invention relates to geometric navigation of high-dimensional
embedding spaces, and more particularly to a method and system for augmenting
pre-computed embedding vectors with a scaled temporal coordinate to create a
unified $(D+1)$-dimensional space in which time is a first-class navigable
geometric axis, and for navigating that space using multi-modal manifold flight
primitives that support continuous interpolation between semantic and temporal
traversal of any timestamped corpus.

% =============================================================================
\section*{BACKGROUND OF THE INVENTION}
\addcontentsline{toc}{section}{Background of the Invention}
% =============================================================================

Standard sentence-transformer models produce fixed-dimensional real-valued
vectors that encode semantic meaning. Two documents discussing similar topics
will land close together in embedding space regardless of when they were
written. A diary entry from 1660 and one from 1669 may occupy nearly identical
coordinates if their semantic content is sufficiently similar. The temporal
dimension is discarded entirely.

Prior art approaches to time-aware embedding include: (a) training specialized
temporal language models from scratch, which requires large datasets and
significant compute; (b) fine-tuning existing models with temporal supervision
signals, which alters the semantic structure of the embedding space; (c)
treating time as a metadata filter applied after nearest-neighbor retrieval,
which does not expose time as a navigable geometric direction; and (d)
concatenating raw Unix timestamps or year integers directly to embedding
vectors, which produces gross magnitude imbalance because raw time values
($\sim 10^9$ for Unix epoch) dwarf normalized embedding magnitudes
($\sim \mathcal{O}(1)$), causing the temporal dimension to dominate retrieval
entirely.

None of the prior art provides: a principled, magnitude-preserving method for
appending a scalar temporal coordinate as a $(D+1)$-th embedding dimension; a
configurable weight parameter that continuously interpolates from
temporally-neutral to temporally-dominant embedding spaces; or a manifold
navigator equipped with time-specific primitives that can face and fly along the
temporal axis, the semantic axes, or any blend thereof.

The parent application (incorporated herein by reference) describes a general
system that pre-computes 768-dimensional embedding vectors for every chunk in an
unstructured text corpus, writes metadata fields --- including timestamps where
available --- directly to each chunk's structured frontmatter during the corpus
emission stage, and stores the embeddings in an aligned JSON cache. The parent
application's metadata-write mechanism is general: it records any metadata field
that can be parsed directly from the source material without neural network
inference. Timestamp extraction is one instance of this general capability,
applicable to corpora such as personal diaries, clinical notes, legal filings,
meeting transcripts, and correspondence, in which each entry carries an
authoritative date.

The specific contribution of the present invention is to promote the timestamp
field from a post-retrieval filter --- the role it plays in the parent application
--- to a first-class geometric dimension of the embedding space itself. This
promotion is the central technical advance claimed herein: by embedding the
temporal coordinate as a geometrically navigable axis rather than treating it
as a metadata attribute to be filtered after retrieval, the present invention
enables continuous, direction-sensitive manifold navigation through time that
is impossible with filter-based approaches. The present invention builds upon
the parent's cache and metadata-write output to create this navigable
temporally-augmented space.

% =============================================================================
\section*{SUMMARY OF THE INVENTION}
\addcontentsline{toc}{section}{Summary of the Invention}
% =============================================================================

The present invention provides a method and system for temporal dimension
augmentation of embedding vectors and multi-modal manifold navigation of the
resulting augmented space. While the parent application applies to any
unstructured text corpus, the present invention specifically exploits temporal
metadata present in one class of corpora --- personal diaries, clinical notes,
legal filings, meeting transcripts, and correspondence --- in which each entry
carries an authoritative timestamp that the parent's direct-metadata-write
mechanism records in chunk frontmatter. The present invention promotes that
timestamp from a retrieval filter to a first-class navigable geometric axis in
the embedding space produced by the parent method.

\textbf{The augmentation method} appends a scalar temporal coordinate to each
$D$-dimensional embedding vector to produce a $(D+1)$-dimensional vector. The
temporal coordinate is first z-score normalized across the corpus, then scaled
to match the mean $L_2$ norm of the embedding vectors, multiplied by a
configurable weight parameter \alphabar. This produces a unified embedding space
in which the temporal axis has a magnitude contribution commensurate with a
single typical embedding axis when $\alpha = 1$, smaller when $\alpha < 1$, and
larger when $\alpha > 1$.

\textbf{The navigation system} builds a $k$-nearest-neighbor graph over the
augmented $(D+1)$-dimensional space and wraps the \turtlend\ N-dimensional
manifold navigator from the \texttt{proteusPy} library with three new time
primitives: \expanddim\ to grow the navigator's basis by one dimension;
\orientintime\ to rotate the navigator's heading to face the temporal axis; and
\orienttoward\ to rotate the heading toward an arbitrary direction vector.

\textbf{Three flight modes} are defined: \emph{semantic flight}, in which the
navigator steers toward the semantically most-distant reachable node; \emph{temporal
flight}, in which the heading is oriented along the temporal axis and the
navigator advances through purely chronological order; and \emph{mixed flight},
in which the heading blends semantic and temporal directions via a configurable
\texttt{time\_blend} parameter on the interval $[0.0,\, 1.0]$.

\textbf{The generalization principle}: the augmentation and navigation machinery
is not specific to time. Any scalar quantity --- sentiment score, reading level,
geographic coordinate, market volatility, author index, topic cluster centroid
--- can be appended as an extra navigable axis using the same algorithm and the
same primitives.

The invention has been reduced to practice and validated on a 8,413-entry
corpus producing 769-dimensional augmented embeddings (768 semantic + 1
temporal) with a KNN graph at $k = 10$.

% =============================================================================
\section{DETAILED DESCRIPTION OF THE PREFERRED EMBODIMENTS}
% =============================================================================

\subsection{System Overview and Relationship to Parent Application}

The present invention receives as input the JSON embedding cache produced by
the multi-process parallel embedder described in the parent application. That
cache contains three index-aligned arrays:

\begin{itemize}[noitemsep]
  \item \texttt{embeddings}: an $N \times D$ float32 array of pre-computed
    sentence-transformer embeddings, where $N$ is the number of corpus chunks
    and $D = 768$ in the preferred embodiment using the
    \texttt{all-mpnet-base-v2} model;
  \item \texttt{texts}: a length-$N$ array of the corresponding chunk text
    strings;
  \item \texttt{timestamps}: a length-$N$ array of ISO-format timestamp strings
    written directly from source metadata during the corpus emission stage
    described in the parent application.
\end{itemize}

The present invention converts the ISO timestamps to floating-point values
(e.g., fractional years or Unix epoch seconds), augments the $N \times D$
embedding matrix to $N \times (D+1)$ using the \augtime\ algorithm described
below, constructs a $k$-nearest-neighbor graph over the augmented space, and
provides the \temporalflyer\ navigator for multi-modal traversal.

No modifications to the parent application's pipeline are required. The
augmentation is applied entirely in memory, post-cache, without modifying the
stored embeddings.

% --------------------------------------------------------------------------
\subsection{Configurable Temporal Dimension Augmentation}
\label{sec:augment}
% --------------------------------------------------------------------------

\subsubsection{The Problem of Magnitude Imbalance}

A 768-dimensional sentence-transformer embedding vector has a mean $L_2$ norm
of approximately 3--8 units in normalized practice. A raw timestamp expressed
as, for example, a Unix epoch integer has a magnitude on the order of
$10^9$\,--\,$10^{10}$. Appending this raw value as a 769th dimension would
cause the temporal component to dominate all nearest-neighbor computations,
collapsing the semantic structure of the space. Conversely, expressing
timestamps as year integers (1660--1669 for the Pepys corpus) still produces
values orders of magnitude different from the embedding norm without scaling.

The invention solves this problem through magnitude-preserving normalization.

\subsubsection{The \augtime\ Algorithm}

Let $\mathbf{E} \in \mathbb{R}^{N \times D}$ denote the embedding matrix and
let $\mathbf{t} \in \mathbb{R}^N$ denote the vector of scalar time values
(e.g., fractional year, Unix epoch, or any monotone temporal encoding).

The algorithm proceeds in three steps:

\textbf{Step 1 --- Z-score normalization of time values.}
\begin{equation}
  \tilde{t}_i = \frac{t_i - \bar{t}}{\sigma_t}
  \label{eq:zscore}
\end{equation}
where $\bar{t}$ is the mean and $\sigma_t$ is the standard deviation of the
time values across the corpus. This centers and scales the temporal coordinate
to unit variance, removing dependence on the choice of time encoding units.

\textbf{Step 2 --- Scaling to match embedding magnitude.}
\begin{equation}
  \hat{t}_i = \tilde{t}_i \cdot \alpha \cdot \frac{\bar{\|\mathbf{e}\|}}{\sqrt{D}}
  \label{eq:scale}
\end{equation}
where $\bar{\|\mathbf{e}\|}$ is the mean $L_2$ norm of the embedding rows
(i.e., $\frac{1}{N}\sum_{i=1}^{N} \|\mathbf{e}_i\|_2$), $D$ is the embedding
dimension, and \alphabar\ is the configurable temporal weight parameter.

The factor $\bar{\|\mathbf{e}\|}/\sqrt{D}$ represents the expected per-axis
contribution of a single embedding dimension. Multiplying by this factor ensures
that $\hat{t}$ has the same magnitude as one typical embedding axis when
$\alpha = 1$.

\textbf{Step 3 --- Column concatenation.}
\begin{equation}
  \mathbf{Z} = \bigl[\,\mathbf{E} \mid \hat{\mathbf{t}}\,\bigr]
  \in \mathbb{R}^{N \times (D+1)}
  \label{eq:augment}
\end{equation}
The augmented matrix $\mathbf{Z}$ replaces $\mathbf{E}$ for all downstream
computations including KNN graph construction and manifold navigation.

\subsubsection{The Alpha Parameter}

The weight parameter $\alpha \geq 0$ provides continuous control over the
relative influence of the temporal dimension:

\begin{itemize}[noitemsep]
  \item $\alpha = 0$: the temporal column is zeroed; $\mathbf{Z}$ reduces to
    $\mathbf{E}$ with an appended zero; temporal ordering has no influence on
    navigation.
  \item $\alpha = 1$: the temporal dimension contributes as much as one
    typical semantic axis; semantic and temporal structure are balanced.
  \item $\alpha > 1$: the temporal dimension is amplified above the average
    semantic axis; temporal proximity begins to dominate nearest-neighbor
    structure.
\end{itemize}

In the reference implementation, $\alpha = 1.0$ was used for all flight
experiments described in Section~\ref{sec:results}.

\subsubsection{Reference Implementation}

Figure~\ref{fig:augment} shows the complete reference implementation of the
\augtime\ function.

\begin{figure}[H]
\begin{lstlisting}[
  caption={Reference implementation of \texttt{augment\_with\_time} --- temporal dimension augmentation with magnitude-preserving scaling.},
  label={fig:augment}
]
import numpy as np
from math import sqrt


def augment_with_time(
    embeddings: np.ndarray,
    time_values: np.ndarray,
    alpha: float = 1.0,
) -> np.ndarray:
    """Append a scaled temporal coordinate as a (D+1)-th dimension.

    Args:
        embeddings:   (N, D) float32 array of embedding vectors.
        time_values:  (N,) float array of scalar time values in any
                      monotone encoding (fractional year, Unix epoch,
                      ordinal day, etc.).
        alpha:        Temporal weight parameter. alpha=0 recovers the
                      original embedding space; alpha=1 makes the
                      temporal axis contribute as much as one typical
                      embedding dimension; alpha>1 makes time dominant.

    Returns:
        (N, D+1) float32 array with the scaled temporal coordinate
        appended as the final column.
    """
    # Step 1: z-score normalize the time values across the corpus
    t_mean = time_values.mean()
    t_std  = time_values.std()
    t_norm = (time_values - t_mean) / t_std          # zero mean, unit var

    # Step 2: scale to match the per-axis embedding magnitude
    # mean_norm / sqrt(D) is the expected contribution of one axis
    mean_norm = np.linalg.norm(embeddings, axis=1).mean()
    D         = embeddings.shape[1]
    t_scaled  = t_norm * alpha * (mean_norm / sqrt(D))

    # Step 3: append the temporal coordinate as the (D+1)-th column
    return np.column_stack([embeddings, t_scaled])
\end{lstlisting}
\end{figure}

% --------------------------------------------------------------------------
\subsection{KNN Graph Construction over Augmented Space}
% --------------------------------------------------------------------------

Following augmentation, the invention constructs a $k$-nearest-neighbor graph
over the rows of $\mathbf{Z}$. In the preferred embodiment, $k = 10$ and
cosine distance is used as the similarity metric. Each node $i$ in the graph
has edges to its $k$ nearest neighbors in the $(D+1)$-dimensional augmented
space.

Because the temporal coordinate is part of the ambient space during KNN
construction, two nodes that are semantically similar but temporally distant
will be slightly farther apart in augmented space than in the original $D$-dimensional
space, and two nodes that are semantically moderate but temporally close will
be slightly nearer. The degree of this effect is controlled by \alphabar.

The KNN graph is the sole data structure consumed by the \temporalflyer\
navigator. It does not require rebuilding when \alphabar\ is changed, provided
the graph was constructed at the desired \alphabar\ value.

% --------------------------------------------------------------------------
\subsection{Time Navigation Primitives}
\label{sec:primitives}
% --------------------------------------------------------------------------

The invention extends the \turtlend\ N-dimensional manifold navigator from the
\texttt{proteusPy} library with three new primitives. \turtlend\ models a
navigator in $\mathbb{R}^N$ as a position vector and an orthonormal frame
(heading, left, up, and higher-dimensional analogues). It advances by stepping
to adjacent nodes in the KNN graph.

\subsubsection{\expanddim\ --- Grow the Basis by One Dimension}

\textbf{Problem:} \turtlend\ is initialized with an $N$-dimensional basis.
When the embedding space is augmented from $D$ to $D+1$ dimensions, the
navigator's internal basis must be extended to accommodate the new axis.

\textbf{Solution:} \expanddim\ appends a new standard basis vector $\mathbf{e}_{D+1}$
(a zero vector with a 1 in the $(D+1)$-th position) to the navigator's
orthonormal frame. The existing $D$ basis vectors are unchanged. No
re-orthogonalization is required because $\mathbf{e}_{D+1}$ is by construction
orthogonal to all vectors in the original $D$-dimensional subspace.

This primitive enables the navigator to be initialized on the original
$D$-dimensional space and then promoted to the augmented $(D+1)$-dimensional
space without rebuilding the navigator from scratch.

\subsubsection{\orientintime\ --- Rotate Heading to Face the Temporal Axis}

\textbf{Problem:} After \expanddim, the navigator knows the temporal axis
exists, but its heading vector points in an arbitrary semantic direction. To
perform temporal flight, the navigator must face forward in time.

\textbf{Solution:} \orientintime\ sets the navigator's heading to
$\mathbf{e}_{D+1}$ (the positive temporal axis direction) or its negation
$-\mathbf{e}_{D+1}$ (backward in time). The remaining basis vectors are
re-orthogonalized via Gram-Schmidt to maintain the orthonormal frame invariant.

After calling \orientintime, stepping forward in the KNN graph advances the
navigator preferentially toward nodes with later timestamps, because the KNN
graph was constructed in the augmented space where temporal proximity is
encoded in the metric.

\subsubsection{\orienttoward\ --- Rotate Heading Toward an Arbitrary Direction}

\textbf{Problem:} Semantic flight requires the navigator to face the
semantically most-distant reachable node, while mixed flight requires a
heading that blends temporal and semantic directions. Neither direction is a
standard basis vector.

\textbf{Solution:} \orienttoward\ accepts an arbitrary direction vector
$\mathbf{v} \in \mathbb{R}^{D+1}$ and rotates the navigator's heading toward
$\hat{\mathbf{v}} = \mathbf{v}/\|\mathbf{v}\|$ via a Householder reflection
that maps the current heading to $\hat{\mathbf{v}}$ while preserving the
orthonormality of the remaining basis vectors. The remaining frame vectors
are updated by the same reflection.

This primitive is the common implementation substrate for both semantic and
mixed flight modes.

% --------------------------------------------------------------------------
\subsection{The TemporalFlyer: Three Flight Modes}
\label{sec:flyer}
% --------------------------------------------------------------------------

The \temporalflyer\ class composes the augmented embedding matrix, the KNN
graph, and the navigation primitives into a complete manifold navigation
system. It maintains internal state including the current position index in
the corpus, the navigator's heading, and the flight log (sequence of visited
node indices and their timestamps).

\subsubsection{Semantic Flight}

In semantic flight, the navigator seeks the semantically most-distant entry at
each step:

\begin{enumerate}[noitemsep]
  \item From the current node $i$, identify the $k$ KNN neighbors.
  \item Compute the augmented-space direction vector from the current position
    to each neighbor.
  \item Select the neighbor whose direction vector has the smallest cosine
    similarity with the current heading (i.e., is most orthogonal or opposed),
    weighted by distance. This steers the navigator toward unexplored semantic
    territory.
  \item Call \orienttoward\ with the selected direction, then advance to the
    selected neighbor.
\end{enumerate}

In semantic flight, time is incidental. The navigator follows meaning; temporal
progression emerges only to the extent that semantically distant entries happen
to be temporally distant. In the reference experiment, semantic flight produced
a Kendall $\tau = +0.19$, indicating weak forward temporal drift consistent
with a corpus whose semantic content changes gradually over time.

\subsubsection{Temporal Flight}

In temporal flight, the navigator orients along the temporal axis and follows
it as directly as the KNN graph permits:

\begin{enumerate}[noitemsep]
  \item Call \orientintime\ to set the heading to $\mathbf{e}_{D+1}$.
  \item At each step, select the KNN neighbor with the maximum projection onto
    the current heading direction (i.e., the neighbor most forward in
    augmented-space time direction).
  \item Advance to the selected neighbor without changing the heading.
\end{enumerate}

Because the KNN graph was constructed in the augmented space, the neighbors
most forward in the temporal direction are nodes that are both temporally
close and semantically proximate. The navigator does not simply traverse
entries in timestamp order; it follows the manifold's geometry, which is
curved relative to the temporal axis. This produces shorter mean time steps
(0.61 years per step in the reference experiment) than semantic flight
(1.56 years per step), with a Kendall $\tau = -0.38$ indicating that the path
is locally coherent in time but globally non-monotone, as expected for a
curved manifold.

\subsubsection{Mixed Flight}

Mixed flight continuously interpolates between semantic and temporal directions
using a configurable \texttt{time\_blend} parameter $\beta \in [0.0,\, 1.0]$:

\begin{enumerate}[noitemsep]
  \item Compute the semantic direction $\mathbf{d}_s$ (direction to the most
    semantically interesting neighbor, as in semantic flight).
  \item Compute the temporal direction $\mathbf{d}_t = \mathbf{e}_{D+1}$.
  \item Compute the blended heading:
    \begin{equation}
      \mathbf{d}_{\text{mixed}} = \frac{(1-\beta)\,\mathbf{d}_s + \beta\,\mathbf{d}_t}
                                       {\|(1-\beta)\,\mathbf{d}_s + \beta\,\mathbf{d}_t\|}
      \label{eq:blend}
    \end{equation}
  \item Call \orienttoward($\mathbf{d}_{\text{mixed}}$) and advance to the
    neighbor with the maximum projection onto the blended heading.
\end{enumerate}

At $\beta = 0$, mixed flight is identical to semantic flight. At $\beta = 1$,
it is identical to temporal flight. Intermediate values produce paths that
spiral through both time and meaning simultaneously, covering broader temporal
spans (2.78 years total in the reference experiment at $\beta = 0.5$) with
intermediate time coherence (Kendall $\tau = -0.15$).

% --------------------------------------------------------------------------
\subsection{Measured Results}
\label{sec:results}
% --------------------------------------------------------------------------

\subsubsection{Experimental Configuration}

All results were obtained on the Pepys corpus as prepared by the parent
application's pipeline.

\begin{center}
\begin{tabular}{ll}
  \toprule
  \textbf{Parameter} & \textbf{Value} \\
  \midrule
  Corpus entries             & 8,413 \\
  Embedding model            & \texttt{all-mpnet-base-v2} \\
  Semantic embedding dim.    & 768 \\
  Augmented dim.\ ($D+1$)    & 769 \\
  Temporal weight $\alpha$   & 1.0 \\
  KNN graph degree $k$       & 10 \\
  Similarity metric          & Cosine \\
  Hardware                   & Apple Silicon (CPU only, no GPU) \\
  \bottomrule
\end{tabular}
\end{center}

\vspace{1em}

Flight experiment parameters:

\begin{center}
\begin{tabular}{ll}
  \toprule
  \textbf{Parameter} & \textbf{Value} \\
  \midrule
  Origin timestamp     & 1663-10-21 \\
  Destination timestamp & 1664-01-23 \\
  Maximum steps        & 150 \\
  \bottomrule
\end{tabular}
\end{center}

\subsubsection{Flight Results}

\begin{center}
\begin{tabular}{lrrrrr}
  \toprule
  \textbf{Mode} & \textbf{Path Length} & \textbf{Monotonicity} &
  \textbf{Kendall $\tau$} & \textbf{Mean $\Delta t$ (yr)} & \textbf{Total Span (yr)} \\
  \midrule
  Semantic                 & 79 steps  & 0.54 & $+0.19$ & 1.56 & 2.01 \\
  Temporal                 & 151 steps & 0.52 & $-0.38$ & 0.61 & 2.36 \\
  Mixed ($\beta = 0.5$)    & 142 steps & 0.52 & $-0.15$ & 1.13 & 2.78 \\
  \bottomrule
\end{tabular}
\end{center}

\textbf{Interpretation.}

\begin{itemize}
  \item \textbf{Semantic flight} ($\tau = +0.19$): weakly forward in time.
    The navigator follows meaning; temporal progression is incidental. The
    corpus's semantic structure drifts slightly forward over the 9-year span,
    explaining the positive $\tau$.

  \item \textbf{Temporal flight} ($\tau = -0.38$): shorter time steps (0.61
    yr/step) and a longer path (151 steps). The negative $\tau$ indicates
    the navigator is following the manifold's geometric curvature, not a
    monotone timestamp sequence. The path stays temporally local at each hop
    but traverses the manifold in a non-monotone global order, confirming
    that the temporal axis is genuinely orthogonal to some semantic axes in
    the augmented space.

  \item \textbf{Mixed flight} ($\beta = 0.5$, $\tau = -0.15$): intermediate
    behavior. Longer mean time steps than temporal (1.13 vs.\ 0.61 yr/step),
    broader total temporal coverage (2.78 yr), and moderate time coherence.
    The navigator balances exploration of semantic space with temporal
    progression.
\end{itemize}

\subsubsection{Manifold Geometry Context}

The reference corpus occupies a 768-dimensional ambient space but has
TwoNN intrinsic dimensionality of 14.26 (reported in the parent application).
The effective manifold is low-dimensional ($\sim 14$\,D) and curved. This
curvature explains why temporal flight produces a non-monotone Kendall $\tau$:
the manifold's geometry does not align with the global time axis, confirming
that the temporal augmentation reveals genuine geometric structure rather than
trivially reordering the corpus by timestamp.

% --------------------------------------------------------------------------
\subsection{Generalization to Arbitrary Scalar Dimensions}
\label{sec:generalization}
% --------------------------------------------------------------------------

The \augtime\ algorithm, the \expanddim\ primitive, the \orienttoward\
primitive, and the three flight modes are not specific to time. The same
machinery applies to any scalar quantity that can be computed for each corpus
entry and carries meaningful ordering or magnitude information.

\textbf{Examples of directly applicable scalar dimensions:}

\begin{center}
\begin{tabular}{ll}
  \toprule
  \textbf{Scalar Quantity} & \textbf{Navigation Capability} \\
  \midrule
  Sentiment score          & Navigate from negative to positive affect \\
  Reading level (Flesch-Kincaid) & Navigate from simple to complex prose \\
  Geographic latitude/longitude & Navigate through spatial distribution \\
  Market volatility index  & Navigate from stable to turbulent periods \\
  Document age             & Navigate from recent to archival content \\
  Author index             & Navigate between author styles \\
  Topic cluster centroid distance & Navigate toward/away from a topic cluster \\
  Citation count           & Navigate from obscure to highly-cited works \\
  \bottomrule
\end{tabular}
\end{center}

In each case, the procedure is identical to the temporal case: compute the
scalar values, call \augtime\ with appropriate \alphabar, build or rebuild the
KNN graph, initialize a \temporalflyer\ (or analogous navigator), and use
\orienttoward\ with the $(D+1)$-th standard basis vector as the target
direction.

Multiple scalar dimensions can be appended simultaneously by applying
\augtime\ iteratively. Each application increases the embedding dimension by
one and adds an independent navigable axis. The navigation primitives
generalize to $D+k$ dimensions for $k$ appended scalars.

Figure~\ref{fig:generalized} shows the generalized call pattern for
multi-scalar augmentation.

\begin{figure}[H]
\begin{lstlisting}[
  caption={Generalized scalar-dimension augmentation --- iterative application to append multiple navigable axes.},
  label={fig:generalized}
]
def augment_with_scalar(
    embeddings: np.ndarray,
    scalar_values: np.ndarray,
    alpha: float = 1.0,
) -> np.ndarray:
    """Append any scalar quantity as a navigable (D+1)-th dimension.

    Identical to augment_with_time() but accepts arbitrary scalar values.
    May be applied iteratively to append multiple axes.

    Example --- append sentiment and time simultaneously:

        Z = augment_with_scalar(embeddings, sentiment_scores, alpha=0.5)
        Z = augment_with_scalar(Z, time_values, alpha=1.0)
        # Z is now (N, D+2): semantic + sentiment + temporal axes

    Args:
        embeddings:     (N, D) float32 embedding matrix.
        scalar_values:  (N,) array of scalar values in any units.
        alpha:          Weight parameter (0=no effect, 1=one axis,
                        >1=amplified).

    Returns:
        (N, D+1) float32 augmented matrix.
    """
    s_mean   = scalar_values.mean()
    s_std    = scalar_values.std()
    s_norm   = (scalar_values - s_mean) / s_std

    mean_norm = np.linalg.norm(embeddings, axis=1).mean()
    D         = embeddings.shape[1]
    s_scaled  = s_norm * alpha * (mean_norm / sqrt(D))

    return np.column_stack([embeddings, s_scaled])
\end{lstlisting}
\end{figure}

% --------------------------------------------------------------------------
\subsection{System Architecture Summary}
% --------------------------------------------------------------------------

Figure~\ref{fig:pipeline} summarizes the complete system, including the parent
application's pipeline stages (shown in boxes with dashed borders) and the
present invention's additions (shown in boxes with solid borders).

\begin{figure}[H]
\begin{lstlisting}[
  language={},
  basicstyle=\ttfamily\small,
  backgroundcolor=\color{codebg},
  frame=single,
  numbers=none,
  caption={Full pipeline including parent application stages and CIP additions.},
  label={fig:pipeline}
]
  PARENT APPLICATION STAGES (incorporated by reference)
  -------------------------------------------------------
  Stage 1: DiaryTransformer -- NLP enrichment, chunking,
           topic/context classification
  Stage 2: Corpus emission -- .md chunks with YAML frontmatter,
           direct temporal provenance write
  Stage 3: DocKG SQLite build -- structural graph, BM25
  Stage 4: DocKG LanceDB build -- 768-dim vector index
  Stage 5: Multi-process Embedder -- parallel sentence-transformer
  Stage 6: JSON Cache -- aligned {embeddings, texts, timestamps}

  PRESENT INVENTION (CIP additions)
  -------------------------------------------------------
  Step A:  Parse timestamps from JSON cache -> float time_values
  Step B:  augment_with_time(embeddings, time_values, alpha)
           -> Z: (N, 769) augmented matrix
  Step C:  Build KNN graph over Z (k=10, cosine distance)
  Step D:  Initialize TemporalFlyer with KNN graph and Z
           - expand_dim()       : extend TurtleND basis to D+1
           - orient_in_time()   : face the temporal axis
           - orient_toward(v)   : face arbitrary direction
  Step E:  Execute flight (semantic | temporal | mixed)
           -> path: [(index, timestamp, text), ...]
\end{lstlisting}
\end{figure}

% =============================================================================
\section{CLAIMS}
% =============================================================================

\noindent\textit{What is claimed is:}

\bigskip

% --------------------------------------------------------------------------
% Independent Claims
% --------------------------------------------------------------------------

\noindent\textbf{INDEPENDENT CLAIMS}

\bigskip

\begin{claim}
A computer-implemented method for augmenting pre-computed embedding vectors
with a temporal coordinate, comprising:

\begin{enumerate}[label=(\alph*), noitemsep]
  \item receiving a matrix $\mathbf{E} \in \mathbb{R}^{N \times D}$ of
    pre-computed floating-point embedding vectors, each row corresponding to
    a text chunk from a timestamped corpus, wherein the embedding vectors
    were computed by a pre-trained sentence-transformer model;

  \item receiving a vector $\mathbf{t} \in \mathbb{R}^N$ of scalar time
    values corresponding index-for-index to the rows of $\mathbf{E}$;

  \item computing z-score normalized time values:
    $\tilde{t}_i = (t_i - \bar{t})\,/\,\sigma_t$,
    where $\bar{t}$ is the mean and $\sigma_t$ is the standard deviation
    of $\mathbf{t}$;

  \item computing a scaled temporal coordinate:
    $\hat{t}_i = \tilde{t}_i \cdot \alpha \cdot \bar{\|\mathbf{e}\|}\,/\,\sqrt{D}$,
    where $\bar{\|\mathbf{e}\|}$ is the mean $L_2$ row norm of $\mathbf{E}$,
    $D$ is the embedding dimension, and $\alpha$ is a configurable weight
    parameter;

  \item forming the augmented matrix:
    $\mathbf{Z} = [\,\mathbf{E}\;\mid\;\hat{\mathbf{t}}\,]
    \in \mathbb{R}^{N \times (D+1)}$
    by appending $\hat{\mathbf{t}}$ as the $(D+1)$-th column;
\end{enumerate}

wherein, when $\alpha = 1$, the magnitude of the appended temporal column is
commensurate with the expected per-axis magnitude of a single dimension of
$\mathbf{E}$, so that temporal and semantic proximity are weighted equally in
subsequent nearest-neighbor computations over $\mathbf{Z}$.
\end{claim}

\begin{claim}
The method of Claim~1, further comprising:

\begin{enumerate}[label=(\alph*), noitemsep]
  \item constructing a $k$-nearest-neighbor graph over the rows of the
    augmented matrix $\mathbf{Z}$, wherein each node corresponds to a corpus
    chunk and each edge connects a node to one of its $k$ nearest neighbors
    under cosine distance in the $(D+1)$-dimensional augmented space;

  \item initializing an $N$-dimensional manifold navigator with a position
    vector and orthonormal basis frame;

  \item extending the navigator's basis from $D$ dimensions to $D+1$
    dimensions by appending a new standard basis vector $\mathbf{e}_{D+1}$
    without modifying the existing $D$ basis vectors;

  \item executing a multi-step traversal of the $k$-nearest-neighbor graph
    that comprises at least one of:
    \begin{enumerate}[label=(\roman*), noitemsep]
      \item orienting the navigator's heading to face the $(D+1)$-th axis
        direction $\mathbf{e}_{D+1}$ and advancing preferentially toward
        neighbors with the greatest projection onto that heading, thereby
        performing temporal flight;
      \item orienting the heading toward a semantically distant neighbor
        direction and advancing preferentially toward nodes maximizing
        semantic distance from the current position, thereby performing
        semantic flight;
      \item blending temporal and semantic heading directions using a
        configurable weight parameter $\beta \in [0, 1]$ and advancing
        toward the neighbor with the greatest projection onto the blended
        heading, thereby performing mixed flight;
    \end{enumerate}

  \item recording the sequence of visited nodes and their associated
    timestamps as a flight path.
\end{enumerate}
\end{claim}

\begin{claim}
A computer-implemented system for multi-modal manifold navigation of a
temporally-augmented embedding space, comprising:

\begin{enumerate}[label=(\alph*), noitemsep]
  \item an augmentation module configured to receive a $D$-dimensional
    embedding matrix $\mathbf{E}$ and a vector of scalar time values
    $\mathbf{t}$, apply z-score normalization to $\mathbf{t}$, scale the
    normalized values to match the mean per-axis magnitude of $\mathbf{E}$
    multiplied by a configurable weight parameter $\alpha$, and return the
    $(D+1)$-dimensional augmented matrix $\mathbf{Z}$;

  \item a graph builder module configured to construct a $k$-nearest-neighbor
    graph over the rows of $\mathbf{Z}$ using cosine distance;

  \item a manifold navigator equipped with:
    \begin{enumerate}[label=(\roman*), noitemsep]
      \item a basis extension primitive that grows the navigator's orthonormal
        frame from $D$ to $D+1$ dimensions by appending the standard basis
        vector $\mathbf{e}_{D+1}$;
      \item a temporal orientation primitive that rotates the navigator's
        heading to face $\pm\mathbf{e}_{D+1}$ and re-orthogonalizes the
        remaining frame vectors;
      \item a directed orientation primitive that rotates the navigator's
        heading toward any specified direction vector via Householder
        reflection while preserving frame orthonormality;
    \end{enumerate}

  \item a flight executor configured to execute at least three navigation
    modes --- semantic flight, temporal flight, and mixed flight --- over the
    $k$-nearest-neighbor graph using the manifold navigator's heading and
    orientation primitives;

  \item wherein the system is configured to receive as input the JSON
    embedding cache produced by the offline pre-computation pipeline of
    \parentapp, without modification to that pipeline.
\end{enumerate}
\end{claim}

\begin{claim}
A computer-implemented method for scalar dimension augmentation and navigation
of embedding spaces, comprising:

\begin{enumerate}[label=(\alph*), noitemsep]
  \item receiving a $D$-dimensional embedding matrix $\mathbf{E} \in
    \mathbb{R}^{N \times D}$ and a vector of scalar values $\mathbf{s} \in
    \mathbb{R}^N$, wherein each scalar value encodes a property of the
    corresponding corpus entry other than semantic content;

  \item applying the augmentation method of Claim~1 with $\mathbf{s}$
    substituted for $\mathbf{t}$ to form a $(D+1)$-dimensional augmented
    matrix;

  \item constructing a $k$-nearest-neighbor graph over the augmented matrix;

  \item navigating the $k$-nearest-neighbor graph using the manifold navigator
    of Claim~3 by orienting the heading along the $(D+1)$-th axis direction
    to traverse the corpus in order of increasing scalar value;

  \item wherein the scalar property is selected from the group consisting of:
    sentiment score, reading level score, geographic coordinate, market
    volatility index, document age, author identifier, topic cluster centroid
    distance, citation count, and any other scalar quantity computable at
    ingest time without neural network inference.
\end{enumerate}
\end{claim}

\bigskip

% --------------------------------------------------------------------------
% Dependent Claims
% --------------------------------------------------------------------------

\noindent\textbf{DEPENDENT CLAIMS}

\bigskip

\begin{claim}
The method of Claim~1, wherein $\alpha = 0$ recovers the original
$D$-dimensional embedding space as a degenerate case of the $(D+1)$-dimensional
augmented space in which the temporal column is identically zero, enabling
the augmentation module to be present in the pipeline without effect when
temporal navigation is not desired.
\end{claim}

\begin{claim}
The method of Claim~1, wherein $\alpha > 1$ causes the temporal dimension to
contribute more than a single typical semantic axis, progressively restricting
the KNN graph's neighborhood structure to temporally proximate entries, and
wherein $\alpha = 1$ is the preferred embodiment that balances semantic and
temporal proximity in the KNN metric.
\end{claim}

\begin{claim}
The method of Claim~2, wherein temporal flight produces a Kendall rank
correlation coefficient $\tau < 0$ between the step index along the flight
path and the timestamp at each step, indicating that the navigator is
traversing the curved manifold geometry rather than simply sorting the corpus
by timestamp, and wherein this non-monotone behavior is a measurable
characteristic of the manifold's curvature relative to the temporal axis.
\end{claim}

\begin{claim}
The method of Claim~2, wherein mixed flight with $\beta = 0.5$ produces a
total temporal span greater than either pure semantic flight or pure temporal
flight over the same maximum number of steps, reflecting the navigator's
ability to explore a broader region of the curved manifold by blending
orthogonal geometric directions.
\end{claim}

\begin{claim}
The system of Claim~3, wherein the manifold navigator is implemented as the
\turtlend\ class from the \texttt{proteusPy} library, and wherein the basis
extension primitive, temporal orientation primitive, and directed orientation
primitive are new methods added to \turtlend\ that do not modify any existing
functionality of that class.
\end{claim}

\begin{claim}
The system of Claim~3, wherein the $k$-nearest-neighbor graph is constructed
with $k = 10$ and cosine distance over a 769-dimensional augmented space
comprising 768-dimensional \texttt{all-mpnet-base-v2} embeddings augmented
with one temporal coordinate at $\alpha = 1.0$, for a corpus of approximately
8,413 entries.
\end{claim}

\begin{claim}
The method of Claim~2, further comprising measuring the temporal coherence
of a flight path by computing the Kendall rank correlation coefficient between
the sequence index and the timestamp sequence of the flight path, wherein:
a coefficient near $+1.0$ indicates monotonically forward temporal traversal;
a coefficient near $0$ indicates temporally uncorrelated traversal;
and a coefficient near $-1.0$ indicates highly locally-coherent but globally
non-monotone traversal characteristic of manifold curvature.
\end{claim}

\begin{claim}
The method of Claim~4, wherein the augmentation method of Claim~1 is applied
iteratively with $k$ distinct scalar value vectors to produce a
$(D+k)$-dimensional augmented matrix, each additional column encoding an
independent navigable axis, and wherein the $k$-nearest-neighbor graph over
the $(D+k)$-dimensional space simultaneously encodes all $k$ scalar properties
in the neighborhood structure.
\end{claim}

\begin{claim}
The method of Claim~1, wherein the scalar time values $\mathbf{t}$ are derived
from a metadata field that was written directly to structured chunk frontmatter
during the corpus emission stage of the general pipeline of \parentapp\ without
any language model extraction, the parent pipeline being capable of writing any
metadata field parseable from source material at ingest time, and wherein the
temporal metadata field is one instance of this general metadata-writing
capability, ensuring 100\% temporal reliability regardless of the language
model's size or temporal reasoning capability.
\end{claim}

\begin{claim}
The method of Claim~13, wherein the corpus belongs to a class of timestamped
corpora --- including personal diaries, clinical notes, legal filings, meeting
transcripts, and correspondence --- in which every entry carries an authoritative
date or timestamp parseable at ingest time, so that the parent pipeline's
direct-metadata-write operation produces a timestamp field in every chunk's
frontmatter, and wherein the present invention reads that timestamp field to
construct the scalar time vector $\mathbf{t}$ used in the augmentation of
Claim~1.
\end{claim}

\begin{claim}
The system of Claim~3, wherein the flight executor records, for each step of
each flight path, at least: the corpus index of the visited node; the
timestamp of the visited node; the text content of the visited node; and the
blended heading direction used to select that node, enabling post-hoc analysis
of the navigator's trajectory through the augmented space.
\end{claim}

\begin{claim}
The method of Claim~2, wherein the $k$-nearest-neighbor graph is reused
unchanged across multiple flight executions with different origin nodes,
destination nodes, or flight modes, such that the $\mathcal{O}(N \cdot D)$
graph construction cost is paid once per $(D+1)$-dimensional augmentation
configuration and amortized across all subsequent flights.
\end{claim}

% =============================================================================
\section*{ABSTRACT}
\addcontentsline{toc}{section}{Abstract}
% =============================================================================

A method and system for appending a scalar temporal coordinate to pre-computed
embedding vectors as a configurable $(D+1)$-th dimension, enabling time to
function as a first-class navigable geometric axis in embedding space. The
augmentation algorithm z-score normalizes the temporal values, scales them to
match the mean per-axis magnitude of the embedding vectors via the factor
$\alpha \cdot \bar{\|\mathbf{e}\|}/\sqrt{D}$, and appends the result as the
final column of the embedding matrix. A configurable weight parameter
\alphabar\ continuously controls the relative influence of the temporal
dimension from neutral ($\alpha = 0$) through balanced ($\alpha = 1$) to
temporally dominant ($\alpha > 1$). A $k$-nearest-neighbor graph constructed
over the augmented space serves a \temporalflyer\ manifold navigator equipped
with three time primitives: \expanddim\ (extend the navigator's basis),
\orientintime\ (face the temporal axis), and \orienttoward\ (face an arbitrary
direction). Three flight modes are defined: semantic flight (navigate by
meaning), temporal flight (navigate by time), and mixed flight (blend both via
a \texttt{time\_blend} parameter). Validated on a 8,413-entry corpus in
769-dimensional augmented space ($k=10$, $\alpha=1.0$), the system produces
measurably distinct navigation behaviors: temporal flight achieves shorter mean
time steps (0.61 yr/step) with Kendall $\tau = -0.38$, reflecting manifold
curvature; semantic flight produces longer steps (1.56 yr/step) with
$\tau = +0.19$; and mixed flight (blend~$= 0.5$) achieves the broadest
temporal coverage (2.78 yr total). The augmentation and navigation pattern
generalizes to any scalar quantity (sentiment, reading level, geographic
coordinate, market volatility), enabling arbitrary navigable axes without
architectural changes. The system receives as input the JSON embedding cache
produced by the offline pre-computation pipeline disclosed in the parent
application, requiring no modification to that pipeline.

% =============================================================================
% END OF DOCUMENT
% =============================================================================

\newpage
\section*{DRAWINGS AND FIGURES}
\addcontentsline{toc}{section}{Drawings and Figures}

The following figures are included in this specification:

\begin{itemize}
  \item \textbf{Figure~1} (Listing~\ref{fig:augment}): Reference implementation
    of \augtime\ showing the three-step algorithm: z-score normalization,
    magnitude-preserving scaling, and column concatenation.

  \item \textbf{Figure~2} (Listing~\ref{fig:generalized}): Generalized
    scalar-dimension augmentation function demonstrating iterative application
    for multiple navigable axes.

  \item \textbf{Figure~3} (Listing~\ref{fig:pipeline}): Full pipeline diagram
    showing the relationship between the parent application's six pipeline
    stages and the five new steps of the present invention.
\end{itemize}

\bigskip
\noindent\rule{\textwidth}{0.4pt}

\bigskip

\noindent\textit{Respectfully submitted,}

\vspace{2em}

\noindent\rule{3in}{0.4pt}

\noindent Eric G. Suchanek, PhD

\noindent Inventor

\noindent Flux-Frontiers

\vspace{1em}

\noindent Filing Date: 2026-03-27

\end{document}
